\chapter*{Реферат}
\addcontentsline{toc}{chapter}{Реферат}
{\centering\bfseries\WorkingTitle\par}
\medskip
{\centering\AuthorName\par}
\medskip
\noindent\textbf{Ключевые слова:} вызванный составной потенциал действия (ECAP), стимуляция спинного мозга, коннектом, \textit{Drosophila melanogaster}, нейродинамика, предобучение, перенос представлений.
\medskip
Работа посвящена биофизическому моделированию кривой роста вызванного составного потенциала действия при стимуляции спинного мозга. Рассматривается возможность предобучения начальных параметров временного модуля на коннектомно-ограниченной динамике \textit{Drosophila melanogaster}. Цель исследования --- установить дополнительную ценность такого предобучения для точности моделирования электрического ответа на независимых человеческих данных.

В первой главе сопоставлены наблюдаемые реакции Drosophila, анатомические ограничения сетевых моделей и механизмы формирования ECAP. Рассмотрены методы предобучения, условия совместимости исходных наблюдений и способы независимой проверки моделей. Показано, какие предположения требуют самостоятельной оценки; экспериментальные результаты предлагаемого переноса не представлены.

Планируется разработать модель нейродинамики и алгоритм генерации синтетических временных рядов, определить переносимую часть параметров и связать целевой временной модуль с отдельным биофизическим оператором. Предусматриваются самостоятельная проверка формирования и регистрации человеческого ответа, сравнение с обучением без переноса и простой моделью кривой роста, а также разграничение влияния структуры, динамики и предобучения. Основной единицей оценки будет участник, основным показателем --- ошибка амплитуды кривой роста. Область применимости будет определяться с учётом чувствительности модели и неопределённости оценки различий.

Предполагаемое практическое применение ограничено численным исследованием измерительного канала. ECAP и нейродинамика Drosophila не рассматриваются как прямые измерения субъективной боли; клинические выводы потребуют отдельной проверки.
